% GENERATED FILE. Edit canonical lessons, then run npm run generate:books.
% canonical-source-hash: fnv1a64:e1e4e5462804fe22
% canonical-lessons: LA-C512-gratus, LA-C512-ingratus, LA-C512-meritum, LA-C512-laus, LA-C512-vitupero, LA-R512-review-505-512

\chapter{Gratitude and praise}
\label{ch:la-gratitude-and-praise}
\hlchaptermodality{512}

\begin{chapteropening}
\textbf{By the end of this chapter:} \emph{I can talk about thanks, praise and blame.}

Complete the last lesson of chapter 512: I can talk about thanks, praise and blame.
\end{chapteropening}

\section[grātus, grāta, grātum]{grātus, grāta, grātum — welcome, pleasing; grateful}

\label{lesson:LA-C512-gratus}



\begin{quote}
Before the new one: say the Latin for to ease, then the Latin for tame.
\end{quote}

\subsection*{You'll want to know: grātus, grāta, grātum}
\textbf{grātus, grāta, grātum} — \textquotedblleft{}welcome, pleasing; grateful\textquotedblright{}.

\begin{quote}
Epistula tua mihi grāta est.
\end{quote}

\begin{quote}
\emph{Your letter is welcome to me.}
\end{quote}

Two sides: pleasing to someone (with the dative), and grateful for something. \emph{Grātiās agō}, already known, belongs here. English \emph{grateful}.

\begin{cousinweb}
\textbf{grātus} is inherited and traced to a Proto-Indo-European root meaning \textquotedblleft{}to praise, welcome\textquotedblright{} (EtymDB); Lewis \& Short compares Greek \emph{charis}. English \emph{grateful} and \emph{gratitude} come from its family.
\end{cousinweb}

\begin{grammarlens}[title={What you've built}]
One more word for telling, explaining and describing at B1.
\end{grammarlens}

\subsection*{Your turn}
\begin{itemize}
\raggedright
  \item \emph{Say it:} \emph{grātus, grāta, grātum}
  \item \emph{Say it:} \emph{grātus, grāta, grātum}, once more
  \item \emph{Say it:} a sentence of your own with \emph{grātus, grāta, grātum}
  \item \emph{Recall:} say the Latin for to ease, then the Latin for tame, then say \emph{grātus, grāta, grātum} again
\end{itemize}

\begin{tcolorbox}[breakable,colback=teal!4,colframe=teal!35!black,title={Before you move on}]
What does grātus, grāta, grātum mean? (\textquotedblleft{}Welcome, pleasing; grateful\textquotedblright{}.) Say it once more.
\end{tcolorbox}
\section[ingrātus, ingrāta, ingrātum]{ingrātus, ingrāta, ingrātum — ungrateful; unpleasant}

\label{lesson:LA-C512-ingratus}



\begin{quote}
Before the new one: say the Latin for tame, then the Latin for welcome.
\end{quote}

\subsection*{You'll want to know: ingrātus, ingrāta, ingrātum}
\textbf{ingrātus, ingrāta, ingrātum} — \textquotedblleft{}ungrateful; unpleasant\textquotedblright{}.

\begin{quote}
Fīlius ingrātus parentēs nōn vīsitat.
\end{quote}

\begin{quote}
\emph{The ungrateful son does not visit his parents.}
\end{quote}

\emph{in-} 'not' + \emph{grātus}. English \emph{ingrate}.

\begin{cousinweb}
\textbf{ingrātus} is built on \emph{in-} \textquotedblleft{}not\textquotedblright{} and \emph{grātus} \textquotedblleft{}pleasing, grateful\textquotedblright{} (Lewis \& Short). English \emph{ingratitude} and \emph{ingrate} come from it.
\end{cousinweb}

\begin{grammarlens}[title={What you've built}]
One more word for telling, explaining and describing at B1.
\end{grammarlens}

\subsection*{Your turn}
\begin{itemize}
\raggedright
  \item \emph{Say it:} \emph{ingrātus, ingrāta, ingrātum}
  \item \emph{Say it:} \emph{ingrātus, ingrāta, ingrātum}, once more
  \item \emph{Say it:} a sentence of your own with \emph{ingrātus, ingrāta, ingrātum}
  \item \emph{Recall:} say the Latin for tame, then the Latin for welcome, then say \emph{ingrātus, ingrāta, ingrātum} again
\end{itemize}

\begin{tcolorbox}[breakable,colback=teal!4,colframe=teal!35!black,title={Before you move on}]
What does ingrātus, ingrāta, ingrātum mean? (\textquotedblleft{}Ungrateful; unpleasant\textquotedblright{}.) Say it once more.
\end{tcolorbox}
\section[meritum, meritī]{meritum, meritī — a kindness, service; merit}

\label{lesson:LA-C512-meritum}



\begin{quote}
Before the new one: say the Latin for welcome, then the Latin for ungrateful.
\end{quote}

\subsection*{You'll want to know: meritum, meritī}
\textbf{meritum, meritī} — \textquotedblleft{}a kindness, service; merit\textquotedblright{}.

\begin{quote}
Prō meritīs tuīs tibi grātiās agō.
\end{quote}

\begin{quote}
\emph{I thank you for your kindnesses.}
\end{quote}

Neuter, second declension, from \emph{mereō} — to earn. A good deed that deserves thanks. English \emph{merit}.

\begin{cousinweb}
\textbf{meritum} is the neuter past participle of \emph{mereō} \textquotedblleft{}to earn, deserve\textquotedblright{}, used as a noun (Lewis \& Short; EtymDB). English \emph{merit} comes from it.
\end{cousinweb}

\begin{grammarlens}[title={What you've built}]
One more word for telling, explaining and describing at B1.
\end{grammarlens}

\subsection*{Your turn}
\begin{itemize}
\raggedright
  \item \emph{Say it:} \emph{meritum, meritī}
  \item \emph{Say it:} \emph{meritum, meritī}, once more
  \item \emph{Say it:} a sentence of your own with \emph{meritum, meritī}
  \item \emph{Recall:} say the Latin for welcome, then the Latin for ungrateful, then say \emph{meritum, meritī} again
\end{itemize}

\begin{tcolorbox}[breakable,colback=teal!4,colframe=teal!35!black,title={Before you move on}]
What does meritum, meritī mean? (\textquotedblleft{}A kindness, service; merit\textquotedblright{}.) Say it once more.
\end{tcolorbox}
\section[laus, laudis]{laus, laudis — praise}

\label{lesson:LA-C512-laus}



\begin{quote}
Before the new one: say the Latin for ungrateful, then the Latin for a kindness.
\end{quote}

\subsection*{You'll want to know: laus, laudis}
\textbf{laus, laudis} — \textquotedblleft{}praise\textquotedblright{}.

\begin{quote}
Mīles laudem ducis accipit.
\end{quote}

\begin{quote}
\emph{The soldier receives the general's praise.}
\end{quote}

Feminine, third declension. \emph{Laudō}, already known, comes from it. English \emph{laud}.

\begin{cousinweb}
\textbf{laus} is a native Latin noun, and \emph{laudō} \textquotedblleft{}to praise\textquotedblright{} is built on it (Lewis \& Short). Lewis \& Short links it with \emph{clueō} \textquotedblleft{}to be spoken of\textquotedblright{}, but the origin is not settled. English \emph{laud} and \emph{laudable} come from the family.
\end{cousinweb}

\begin{grammarlens}[title={What you've built}]
One more word for telling, explaining and describing at B1.
\end{grammarlens}

\subsection*{Your turn}
\begin{itemize}
\raggedright
  \item \emph{Say it:} \emph{laus, laudis}
  \item \emph{Say it:} \emph{laus, laudis}, once more
  \item \emph{Say it:} a sentence of your own with \emph{laus, laudis}
  \item \emph{Recall:} say the Latin for ungrateful, then the Latin for a kindness, then say \emph{laus, laudis} again
\end{itemize}

\begin{tcolorbox}[breakable,colback=teal!4,colframe=teal!35!black,title={Before you move on}]
What does laus, laudis mean? (\textquotedblleft{}Praise\textquotedblright{}.) Say it once more.
\end{tcolorbox}
\section[vituperō, vituperāre]{vituperō, vituperāre — to blame, scold}

\label{lesson:LA-C512-vitupero}



\begin{quote}
Before the new one: say the Latin for a kindness, then the Latin for praise.
\end{quote}

\subsection*{You'll want to know: vituperō, vituperāre}
\textbf{vituperō, vituperāre} — \textquotedblleft{}to blame, scold\textquotedblright{}.

\begin{quote}
Magister discipulum pigrum vituperat.
\end{quote}

\begin{quote}
\emph{The teacher scolds the lazy pupil.}
\end{quote}

\emph{vituperō, vituperāre, vituperāvī, vituperātum}, the opposite of \emph{laudō}. English \emph{vituperation}.

\begin{cousinweb}
\textbf{vituperō} is built on \emph{vitium} \textquotedblleft{}a fault\textquotedblright{} and \emph{parō} \textquotedblleft{}to prepare\textquotedblright{} (Lewis \& Short; EtymDB): to make out a fault. English \emph{vituperate} comes from it.
\end{cousinweb}

\begin{grammarlens}[title={What you've built}]
One more word for telling, explaining and describing at B1.
\end{grammarlens}

\subsection*{Your turn}
\begin{itemize}
\raggedright
  \item \emph{Say it:} \emph{vituperō, vituperāre}
  \item \emph{Say it:} \emph{vituperō, vituperāre}, once more
  \item \emph{Say it:} a sentence of your own with \emph{vituperō, vituperāre}
  \item \emph{Recall:} say the Latin for a kindness, then the Latin for praise, then say \emph{vituperō, vituperāre} again
\end{itemize}

\begin{tcolorbox}[breakable,colback=teal!4,colframe=teal!35!black,title={Before you move on}]
What does vituperō, vituperāre mean? (\textquotedblleft{}To blame, scold\textquotedblright{}.) Say it once more.
\end{tcolorbox}
\section[(dialogue)]{Words from chapters 505–512 — a review}

\label{lesson:LA-R512-review-505-512}



\begin{quote}
Say the Latin for merciful and mercy.
\end{quote}

\subsection*{Your turn}
Say these aloud:

\begin{itemize}
\raggedright
  \item merciful, mercy, kindness, kindly, goodwill
  \item courteous, friendliness, refinement, country manners, generosity
  \item pride, arrogant, arrogance, insolent, boastfulness
  \item cowardice, sloth, sluggish, vigorous, diligent
  \item courage, boldness, rash, steadfast, good sense
  \item silly, mad, madness, out of one's mind, to be delirious
  \item cruelty, fury, to hurt, to ease, tame
  \item welcome, ungrateful, a kindness, praise, to blame
\end{itemize}

\begin{tcolorbox}[breakable,colback=teal!4,colframe=teal!35!black,title={Before you move on}]
Merciful? (\textbf{clēmēns, clēmentis}.) Mercy? (\textbf{clēmentia, clēmentiae}.) Kindness? (\textbf{benignitās, benignitātis}.) Kindly? (\textbf{benevolus, benevola, benevolum}.) Goodwill? (\textbf{benevolentia, benevolentiae}.) Courteous? (\textbf{cōmis, cōme}.) Friendliness? (\textbf{cōmitās, cōmitātis}.) Refinement? (\textbf{urbānitās, urbānitātis}.) Country manners? (\textbf{rūsticitās, rūsticitātis}.) Generosity? (\textbf{līberālitās, līberālitātis}.) Pride? (\textbf{superbia, superbiae}.) Arrogant? (\textbf{arrogāns, arrogantis}.) Arrogance? (\textbf{arrogantia, arrogantiae}.) Insolent? (\textbf{īnsolēns, īnsolentis}.) Boastfulness? (\textbf{vānitās, vānitātis}.) Cowardice? (\textbf{ignāvia, ignāviae}.) Sloth? (\textbf{pigritia, pigritiae}.) Sluggish? (\textbf{segnis, segne}.) Vigorous? (\textbf{strēnuus, strēnua, strēnuum}.) Diligent? (\textbf{assiduus, assidua, assiduum}.) Courage? (\textbf{fortitūdō, fortitūdinis}.) Boldness? (\textbf{audācia, audāciae}.) Rash? (\textbf{temerārius, temerāria, temerārium}.) Steadfast? (\textbf{cōnstāns, cōnstantis}.) Good sense? (\textbf{prūdentia, prūdentiae}.) Silly? (\textbf{ineptus, inepta, ineptum}.) Mad? (\textbf{īnsānus, īnsāna, īnsānum}.) Madness? (\textbf{īnsānia, īnsāniae}.) Out of one's mind? (\textbf{dēmēns, dēmentis}.) To be delirious? (\textbf{dēlīrō, dēlīrāre}.) Cruelty? (\textbf{crūdēlitās, crūdēlitātis}.) Fury? (\textbf{saevitia, saevitiae}.) To hurt? (\textbf{laedō, laedere}.) To ease? (\textbf{mītigō, mītigāre}.) Tame? (\textbf{mānsuētus, mānsuēta, mānsuētum}.) Welcome? (\textbf{grātus, grāta, grātum}.) Ungrateful? (\textbf{ingrātus, ingrāta, ingrātum}.) A kindness? (\textbf{meritum, meritī}.) Praise? (\textbf{laus, laudis}.) To blame? (\textbf{vituperō, vituperāre}.)
\end{tcolorbox}
